\chapter{Integration: Jarvis-Operas}
\label{ch:jarvis-operas}

\newthought{\Operas\ is the operator library} behind the \yamlkey{Operas} workflow backend
(Chapter~\ref{ch:operas}). It is a registry of named Python functions that \Jarvis\ can look up and
call. This chapter explains the registry, its \cmd{jopera} command, and how your task card reaches
into it; \Operas\ has its own documentation for the full \textsc{api}.

\section{The registry model}
\label{sec:op-registry}

\Operas\ registers Python callables under stable names and resolves them on demand. Every operator
has a two-part name:

\begin{shellbox}
<namespace>:<name>      # e.g. math:add, stat:chi2_cov, helper:eggbox
\end{shellbox}

\noindent An operator can be called synchronously or asynchronously---which is exactly what an
\yamlkey{Operas} module's \yamlkey{call\_mode} (\code{call} or \code{acall}) selects
(Chapter~\ref{ch:operas}). \Operas\ comes bundled with \Jarvis\ (Chapter~\ref{ch:installation}).

\section{Built-in operators}
\label{sec:op-builtins}

A few operators ship ready to use (Table~\ref{tab:op-builtins})---enough to build and test a
workflow before adding your own.

\begin{table}[ht]
  \footnotesize
  \begin{tabular}{@{}ll@{}}
    \toprule
    \textbf{Operator} & \textbf{Purpose} \\
    \midrule
    \code{math:add(a, b)}            & add two values \\
    \code{math:identity(x)}          & return the input unchanged \\
    \code{stat:chi2\_cov(residual, cov)} & a covariance-matrix $\chi^2$ term \\
    \code{helper:eggbox(inputs)}     & the EggBox helper used by the quick-starts \\
    \bottomrule
  \end{tabular}
  \caption{Selected built-in \Operas\ operators.}
  \label{tab:op-builtins}
\end{table}

\section{Discovering operators with \texttt{jopera}}
\label{sec:op-cli}

The \cmd{jopera} command lets you explore and test the registry from the shell---a fast way to
check an operator's name and signature before wiring it into a card:

\begin{shellbox}
jopera list                       # all operators
jopera list --namespace math      # one namespace
jopera info stat:chi2_cov --json  # signature and metadata
jopera call helper:eggbox --kwargs '{"inputs":{"x":0.5,"y":0.0}}'
\end{shellbox}

\noindent \cmd{jopera} also has \code{acall} (async), \code{load} (load user operators), and
\code{help}.

\section{How a task card uses it}
\label{sec:op-card}

Your card reaches the registry in two ways, both from Chapter~\ref{ch:operas}:

\begin{description}[style=nextline, leftmargin=1.2em]
  \item[As workflow modules] an \yamlkey{Operas} module names an \yamlkey{operator} and maps its
        \yamlkey{input}/\yamlkey{output}---the operator becomes a step in the workflow.
  \item[As expression functions] the \yamlkey{Operas.Functions} whitelist exposes chosen operators
        to symbolic expressions, so you can call them inside a \yamlkey{LogLikelihood} term
        (Chapter~\ref{ch:symbols}).
\end{description}

\section{Your own operators}
\label{sec:op-user}

When a calculation does not exist as a built-in, write it as a Python callable and load it from a
file---no need to wrap an external program:

\begin{shellbox}
jopera call my_ops:my_score --user-ops /path/to/my_ops.py
\end{shellbox}

\noindent Once registered, your operator is referenced from an \yamlkey{Operas} module just like a
built-in. This is the lightest way to add a custom observable or score to a scan.

\section{Summary}
\label{sec:op-jarvis-summary}

\Operas\ is a named registry (\code{namespace:name}) of Python operators that the \yamlkey{Operas}
backend calls. Explore it with \cmd{jopera}, use built-ins like \code{helper:eggbox} or load your
own, and reference them from \yamlkey{Operas} modules and expressions. For the full registry
\textsc{api} and operator catalog, see the \Operas\ documentation.
